\chapter{结论与展望}\label{chap6_conclusion_outlook}

\section{主要研究结论}\label{sec_ch6_conclusion}

本文面向塔吊吊装作业中的近场安全防护需求，围绕“吊装构件识别、空间定位、风险预警与避障设计”这一主线，构建了从视觉感知到三维安全表达、再到预警和控制接口设计的完整研究框架。与此前仅强调图像检测或单一距离监测的思路相比，本文更加突出吊装构件在三维空间中的位置与边界表达，强调预警判断必须建立在可解释的空间几何关系之上。围绕这一目标，本文的主要研究结论可概括为以下三个方面。

第一，本文完成了基于 YOLO26 的吊装构件视觉感知与边缘端部署路径构建，并明确了视觉侧在系统中的职责边界。研究结果表明，YOLO26 分割模型能够为吊装构件提供区域级语义结果，其价值不在于独立承担全部风险判断，而在于为后续点云投影和空间定位提供可靠的二维约束。通过 TensorRT 与 FP16 半精度部署，视觉感知链路已经具备在 Jetson 边缘平台上在线运行的工程基础，从而为整套系统的实时实现提供了前提条件。

第二，本文构建了基于分割掩码与点云投影回查的吊装构件空间定位方法，打通了从二维语义到三维安全表达的关键环节。研究表明，将 LiDAR 点逐点投影到图像平面，再依据分割掩码回查目标点云集合，是一种兼顾几何可靠性与语义可解释性的融合路径。该方法能够在现有系统中稳定输出吊装构件实例点集、三维质心以及语义着色点云，证明了利用视觉分割结果约束点云空间表达的可行性。进一步地，本文提出以质心、空间包络和边界距离等形式构建预警所需的距离基准，为吊装场景中的安全边界计算提供了统一框架。

第三，本文在已有动态检测与轨迹预测基础上，形成了面向吊装作业的多层级预警与避障设计框架。研究指出，静态距离阈值卡控能够作为风险筛查的基础层，动态点检测与卡尔曼预测则为相向接近、最近接距离估计和 TTC 风险升级提供时间维度依据。基于此，本文进一步提出三级预警状态机、统一预警输出字段以及与控制接口解耦的动作建议机制，使系统研究范围从“感知能否成立”扩展到“风险如何分级、动作如何承接”的完整闭环框架。虽然其中部分内容仍属于论文设计方案，但它们建立在现有代码已具备的感知和预测基础之上，具有明确的工程实现路径。

总体而言，本文的研究表明，吊装安全预警系统不应仅把视觉检测、点云处理和风险判断看作彼此独立的模块，而应围绕吊装构件这一核心对象建立统一的空间语义链路。只有当视觉识别、三维定位、动态趋势分析和控制接口设计在同一框架下被组织起来时，系统才能真正具备从感知走向预警、再走向安全辅助控制的扩展能力。

\section{论文主要创新点}\label{sec_ch6_innovation}

结合全文研究内容，本文的主要创新点可归纳为以下四个方面。

第一，提出了面向吊装构件实时识别任务的 YOLO26 边缘端视觉感知方案。与将视觉任务简单理解为通用目标检测不同，本文将视觉侧明确定位为吊装构件或吊物的区域级识别与分割，并结合 TensorRT 与 FP16 部署形成 Jetson 平台上的在线推理路径。该处理使视觉模型不再是独立的离线精度工具，而成为后续空间定位链路的稳定前端入口。

第二，提出了基于分割掩码约束的点云投影回查三维定位方法。本文不是由图像像素反求唯一三维点，而是充分利用 LiDAR 已具备的直接测距能力，采用“3D 点投影到 2D 图像，再按 mask 回查点云集合”的思路建立视觉与点云的对应关系。该方法将二维语义边界和三维几何点集自然结合，为吊装构件的质心提取、空间包络构建和最小距离计算提供了统一基础。

第三，提出了融合静态距离卡控与动态趋势预测的多层级预警链路设计。本文在现有动态点检测、聚类和卡尔曼预测基础上，进一步将静态距离阈值、相对运动方向、预测最近接距离和 TTC 指标纳入统一风险判级框架，构建了适用于吊装作业场景的三级预警状态机。这一设计使预警逻辑不再停留在单一瞬时距离判断，而能够体现相向接近场景中的时间提前量价值。

第四，提出了面向控制接口级应用的避障策略框架。本文没有把研究止步于感知和风险显示，而是进一步给出了风险状态消息、控制建议消息及分级动作策略的接口组织方法，使系统能够从“输出风险提示”自然扩展到“生成限速、限向或停止建议”的安全辅助控制框架。该创新点的意义在于，它为后续实机控制接口接入预留了完整的数据和逻辑边界。

\section{研究不足与未来展望}\label{sec_ch6_outlook}

尽管本文已经完成了从吊装构件视觉识别到空间定位、再到预警与避障设计的整体框架构建，但仍存在若干需要进一步完善之处。

第一，视觉模型的训练细节与定量精度结果尚未完全整理。当前研究已确认数据来源于工地实采样本，也已经完成 YOLO26 的预测与部署链路，但关于类别定义、训练平台、关键超参数以及分组测试指标等信息仍需进一步系统化整理。这意味着本文在训练与验证部分仍保留若干 \textbf{【待补充】} 项，后续工作应围绕数据集标注口径、实验配置和结果统计建立更完整的实验档案，以增强论文在模型训练层面的完备性。

第二，基于空间距离的多层级预警和相向运动风险升级虽然已经形成明确设计，但尚未在当前仓库中实现完整闭环。当前系统已经具备吊装构件三维定位、动态点检测和卡尔曼预测基础，这为 TTC 计算、未来最近接距离估计和三级状态机实现提供了充分条件；但真正的统一预警消息定义、分层判级逻辑和控制接口仍需进一步工程化开发。因此，后续研究应优先补齐预警模块和控制接口原型，使第4章提出的设计真正转化为可在线验证的闭环系统。

第三，复杂工况下的泛化能力和参数敏感性仍需更大规模实验验证。吊装作业现场存在光照剧烈变化、背景遮挡、吊物大幅摆动、动态目标多样化以及传感器安装扰动等复杂因素，这些因素都会同时影响视觉分割、点云投影和动态趋势预测效果。未来应进一步开展多场景、长时间和多参数组合实验，重点评估同步窗口、下采样比例、距离阈值、预测时域和动作策略参数对系统稳定性与误报漏报特性的影响，从而建立更稳健的参数标定方法。

第四，系统尚未完成与真实控制系统的深度联调，避障动作仍主要停留在接口级和原型级设计验证阶段。未来可在保证安全性的前提下，将预警系统与塔吊控制器、人机交互界面或现场辅助决策平台进一步联动，探索从“风险提示”到“运动约束建议”再到“安全辅助控制”的渐进式应用路径。同时，还可进一步研究多目标联合风险评估、在线自适应阈值调整和更精细的 TTC 建模方法，以提升系统在复杂吊装场景中的适应能力和实用价值。

综上，本文已经为吊装作业中的吊装构件识别、空间定位、风险预警与避障接口设计奠定了较为完整的研究基础。后续工作若能沿着“补齐实验数据、完成预警闭环、接入控制原型、扩展复杂场景验证”的方向持续推进，则有望将本文提出的方法从研究型原型进一步发展为更具工程可用性的吊装安全辅助系统。

% 参考文献请使用主文件中的 biblatex + \printbibliography 自动生成。